\documentclass[11pt]{article}

\usepackage{amsmath,amssymb,amsthm,mathtools}
\usepackage[margin=1in]{geometry}
\usepackage{microtype}
\usepackage[hidelinks]{hyperref}

\newtheorem{theorem}{Theorem}[section]
\newtheorem{lemma}[theorem]{Lemma}
\newtheorem{corollary}[theorem]{Corollary}
\newtheorem{remark}[theorem]{Remark}

\newcommand{\dd}{\,\mathrm d}

\title{Atlas 145: Concentrating Two Triangle Pages on an Edge of a $4$-Cycle}
\author{}
\date{}

\begin{document}

\maketitle

\begin{abstract}
We prove the chromatic-polynomial lower bound for NetworkX Atlas graph 145
on the full interval $p\in[1/2,1]$ and for every graphon.  The graph is a
$4$-cycle with two triangle pages on one edge.  Moving one page to the
adjacent cycle edge produces Atlas 148, whose full-interval theorem is known.
The difference of the two homomorphism densities is exactly one half of a
square under the $4$-cycle homomorphism measure.  This gives a short,
division-free transfer valid on arbitrary probability spaces.
\end{abstract}


%%%%%%%%%%%%%%%%%%%%%%%%%%%%%%%%%%%%%%%%%%%%%%%%%%%%%%%%%%%%%%%%%%%%%%%%%%%%%%%
\section{The two page placements}
%%%%%%%%%%%%%%%%%%%%%%%%%%%%%%%%%%%%%%%%%%%%%%%%%%%%%%%%%%%%%%%%%%%%%%%%%%%%%%%

Let $(\Omega,\mathcal A,\mu)$ be a probability space and let
$W\colon\Omega^2\to[0,1]$ be symmetric and jointly measurable.  Put
\[
 p=\int_{\Omega^2}W(x,y)\dd\mu^2,
 \qquad
 S(x,y)=\int_\Omega W(x,z)W(y,z)\dd\mu(z).
\]

On labelled vertices $0,1,2,3,4,5$, let $H_{\rm same}$ have edges
\begin{equation}
 01,12,23,30,04,14,05,15.
 \label{eq:same-edges}
\end{equation}
Thus $0$--$1$--$2$--$3$--$0$ is a $4$-cycle and vertices $4,5$ are two triangle pages on
the same cycle edge $01$.  This graph is NetworkX Atlas 145, with graph6
code \texttt{Exe\_}.  Let $H_{\rm adj}$ instead have edges
\begin{equation}
 01,12,23,30,04,14,05,35.
 \label{eq:adjacent-edges}
\end{equation}
It has one triangle page on each of the adjacent cycle edges $01$ and $03$.
This is Atlas 148, with graph6 code \texttt{EyUG}.  The exact isomorphism
checks are reproduced by the verifier cited below.

The triangle in $H_{\rm same}$ forces chromatic number at least three.  A
proper three-colouring has colour classes
\[
 \{0,2\},\qquad \{1,3\},\qquad \{4,5\}.
\]
Hence the chromatic number is three and the admissible density interval is
$p\in[1/2,1]$.

To count its proper $r$-colourings, first colour the page-bearing edge $01$.
Each page then has $r-2$ choices.  Conditional on the two distinct endpoint
colours, the length-three path $0$--$3$--$2$--$1$ has
\[
 (r-1)+(r-2)^2=r^2-3r+3
\]
proper internal colourings.  Consequently
\begin{equation}
 \chi_{H_{\rm same}}(r)
 =r(r-1)(r-2)^2(r^2-3r+3).
 \label{eq:chromatic}
\end{equation}
The same polynomial holds for $H_{\rm adj}$, either by the analogous direct
count or by the computation in the Atlas 148 theorem.  Substitution into the
chromatic target gives
\begin{equation}
 \Phi_{H_{\rm same}}(p)
 =(1-p)^6\chi_{H_{\rm same}}\!\left(\frac1{1-p}\right)
 =p(2p-1)^2(3p^2-3p+1).
 \label{eq:target}
\end{equation}


%%%%%%%%%%%%%%%%%%%%%%%%%%%%%%%%%%%%%%%%%%%%%%%%%%%%%%%%%%%%%%%%%%%%%%%%%%%%%%%
\section{The page-concentration square}
%%%%%%%%%%%%%%%%%%%%%%%%%%%%%%%%%%%%%%%%%%%%%%%%%%%%%%%%%%%%%%%%%%%%%%%%%%%%%%%

For $x_0,x_1,x_2,x_3\in\Omega$, define the finite nonnegative measure
\begin{equation}
 \dd\Lambda
 =W(x_0,x_1)W(x_1,x_2)W(x_2,x_3)W(x_3,x_0)
   \dd\mu(x_0)\dd\mu(x_1)\dd\mu(x_2)\dd\mu(x_3).
 \label{eq:cycle-measure}
\end{equation}
No normalization is imposed; its total mass may be zero.  Put
\[
 A=S(x_0,x_1),\qquad B=S(x_0,x_3).
\]

\begin{lemma}[Adjacent-page concentration]
\label{lem:concentration}
For every graphon on every probability space,
\begin{equation}
 t(H_{\rm same},W)-t(H_{\rm adj},W)
 =\frac12\int_{\Omega^4}(A-B)^2\dd\Lambda\geq0.
 \label{eq:square}
\end{equation}
\end{lemma}

\begin{proof}
Integrating the two page vertices in \eqref{eq:same-edges} and
\eqref{eq:adjacent-edges} gives, respectively,
\[
 t(H_{\rm same},W)=\int A^2\dd\Lambda,
 \qquad
 t(H_{\rm adj},W)=\int AB\dd\Lambda.
\]
The reflection of the $4$-cycle that interchanges $x_1$ and $x_3$ while
fixing $x_0,x_2$ preserves \eqref{eq:cycle-measure}.  Therefore
\[
 \int A^2\dd\Lambda=\int B^2\dd\Lambda.
\]
Expanding the square now gives
\[
 \frac12\int(A-B)^2\dd\Lambda
 =\int A^2\dd\Lambda-\int AB\dd\Lambda,
\]
which is \eqref{eq:square}.  Every integrand is bounded and measurable, so
Tonelli--Fubini and the measure-preserving variable permutation are valid on
an arbitrary probability space.  No quotient or positivity assumption on
$t(C_4,W)$ is used.
\end{proof}

The identity also records its precise equality condition: the two rooted
common-neighbour values agree $\Lambda$-almost everywhere.  We do not need to
classify the graphons satisfying that condition.


%%%%%%%%%%%%%%%%%%%%%%%%%%%%%%%%%%%%%%%%%%%%%%%%%%%%%%%%%%%%%%%%%%%%%%%%%%%%%%%
\section{The Atlas 145 bound}
%%%%%%%%%%%%%%%%%%%%%%%%%%%%%%%%%%%%%%%%%%%%%%%%%%%%%%%%%%%%%%%%%%%%%%%%%%%%%%%

We use the following already established result.

\begin{theorem}[Atlas 148 theorem \cite{Atlas148}]
\label{thm:atlas148}
Every graphon of edge density $p\in[1/2,1]$ satisfies
\[
 t(H_{\rm adj},W)
 \geq p(2p-1)^2(3p^2-3p+1).
\]
\end{theorem}

That theorem is an arbitrary-graphon result.  Its low-density half uses the
$K_3$ case of Reiher's clique density theorem \cite{ReiherCliqueDensity}; its
high-density half uses a two-coordinate Hilbert projection and exact rational
Bernstein certificates.  No finite-step or regularity hypothesis is imported
here.

\begin{corollary}[Atlas 145, graph6 \texttt{Exe\_}]
Every graphon $W$ of edge density $p\in[1/2,1]$ satisfies
\[
 \boxed{
 t(H_{\rm same},W)
 \geq p(2p-1)^2(3p^2-3p+1)
 =\Phi_{H_{\rm same}}(p).
 }
\]
\end{corollary}

\begin{proof}
Lemma~\ref{lem:concentration} and Theorem~\ref{thm:atlas148} give
\[
 t(H_{\rm same},W)\geq t(H_{\rm adj},W)
 \geq\Phi_{H_{\rm same}}(p),
\]
where the last target equality follows from \eqref{eq:target}.  At $p=1/2$
the target is zero, and at $p=1$ the graphon is one almost everywhere and
both sides are one, so the endpoint cases require no limiting argument.
\end{proof}


%%%%%%%%%%%%%%%%%%%%%%%%%%%%%%%%%%%%%%%%%%%%%%%%%%%%%%%%%%%%%%%%%%%%%%%%%%%%%%%
\section{Audit and limitations}
%%%%%%%%%%%%%%%%%%%%%%%%%%%%%%%%%%%%%%%%%%%%%%%%%%%%%%%%%%%%%%%%%%%%%%%%%%%%%%%

The comparison is an exact relocation theorem for these two page placements:
two pages on one edge of a $4$-cycle dominate one page on each of two adjacent
edges.  It does not assert arbitrary triangle-page concentration on an
arbitrary base graph; the reflection symmetry of the $4$-cycle is the exact
hypothesis used.  In particular, putting the two pages on opposite cycle
edges gives Atlas 152 (graph6 \texttt{EhNG}), which has an accepted exact
counterexample; the present one-sided concentration inequality cannot prove
a lower bound for that placement.  The argument also does not revive the
false product comparison
\[
 p\,t(H_{\rm same},W)\geq t(K_4-e,W)t(C_4,W)
\]
recorded in the failed-attempt ledger.

The Atlas identifications, graph6 codes, chromatic polynomials, targets,
density decompositions, and exact square identity on 503 rational step
graphons are independently reconstructed by
\begin{verbatim}
python codes/verify_c4_triangle_page_concentration.py
\end{verbatim}

No Lean file is changed.  A future formalization needs the finite-measure
square identity above and the existing Atlas 148 theorem as its only
graph-specific inputs.  The latter still carries the formal dependency on
the sharp triangle-density theorem.

\begin{thebibliography}{2}
\bibitem{Atlas148}
Atlas 148 proof note,
\path{notes/atlas148_paw_bias_hilbert_projection.tex}.

\bibitem{ReiherCliqueDensity}
C.~Reiher,
\emph{The clique density theorem},
Annals of Mathematics \textbf{184} (2016), 683--707;
\url{https://annals.math.princeton.edu/2016/184-3/p01}.
\end{thebibliography}


%%%%%%%%%%%%%%%%%%%%%%%%%%%%%%%%%%%%%%%%%%%%%%%%%%%%%%%%%%%%%%%%%%%%%%%%%%%%%%%
\appendix
\section{What the Lean formalization actually proves}
%%%%%%%%%%%%%%%%%%%%%%%%%%%%%%%%%%%%%%%%%%%%%%%%%%%%%%%%%%%%%%%%%%%%%%%%%%%%%%%

The corollary is formalized in full, on $p\in[1/2,1]$.  The development is
\texttt{lean/Taeyoung/Methods/Atlas145.lean}, the catalogue theorem is
\texttt{satisfiesLowerBound\_145}, and every declaration depends on exactly
\texttt{[propext, Classical.choice, Quot.sound]}.  Theorem~\ref{thm:atlas148}
is no longer quoted from outside: it is
\texttt{Methods.Atlas148.homDensity\_graph148\_bound}, itself formalized, so
the closing remark of Section~4 is superseded on both counts.

\subsection{The square identity becomes one Cauchy--Schwarz}
\label{sec:app145-cs}

Lemma~\ref{lem:concentration} is proved above by reflecting the four-fold
integral \eqref{eq:cycle-measure} in the map $x_1\leftrightarrow x_3$.  A
choice of labelling removes the reflection entirely.

Integrate the two \emph{opposite} corners $x_0,x_2$ of the $4$-cycle
outermost.  The remaining cycle vertices $x_1$ and $x_3$ then appear in
disjoint factors of $\dd\Lambda$, so for fixed $x_0,x_2$ everything
factorizes through the three arm integrals
\[
 P=\int_\Omega W(x_0,u)W(u,x_2)\dd\mu(u),\quad
 Q=\int_\Omega W(x_0,u)W(u,x_2)S(x_0,u)\dd\mu(u),
\]
\[
 R=\int_\Omega W(x_0,u)W(u,x_2)S(x_0,u)^2\dd\mu(u),
\]
in the form
\[
 t(H_{\rm adj},W)=\int\!\!\int Q^2,
 \qquad
 t(H_{\rm same},W)=\int\!\!\int P\cdot R .
\]
The comparison is then the pointwise weighted Cauchy--Schwarz
\[
 Q^2\leq P\cdot R,
\]
that is $(\int A\eta)^2\leq(\int A)(\int A\eta^2)$ at weight
$A(u)=W(x_0,u)W(u,x_2)$ and $\eta(u)=S(x_0,u)$ --- the same inequality this
project uses throughout, and the only inequality in the whole file.  Note that
$P=S(x_0,x_2)$ is the codegree of the opposite corners, so all three arms are
instances of one operator.

The two formulations agree, of course; \eqref{eq:square} is the same
statement with the Cauchy--Schwarz deficit written as a square.  What the
relabelling buys is that no permutation of a four-fold integral has to be
justified, and no measure-preserving change of variables is needed: the
factorization is ordinary Fubini through the peeling the project already
performs for every row.

\subsection{Transporting the Atlas 148 bound}
\label{sec:app145-transport}

Both graphs are laid out on that frame --- vertices $0,1$ the opposite
corners, $2,3$ the other two cycle vertices, $4,5$ the pages --- so the Lean
copy of Atlas 148 used here is
\[
 02,\;03,\;04,\;05,\;12,\;13,\;24,\;35,
\]
which is \texttt{Methods.Atlas148.graph148} relabelled.  The two densities are
identified by \texttt{Foundation.homDensity\_iso}, so the Atlas 148 theorem
applies to the frame copy without being reproved.

\subsection{Chromatic data}
\label{sec:app145-chromatic}

The polynomial \eqref{eq:chromatic} is proved, not assumed.  Its factor
$r^2-3r+3$ is irreducible over the rationals, so no clique-attachment tower
produces it, and the route is the same as for Atlas 148: the six surjective
colouring counts
\[
 0,\;0,\;0,\;18,\;264,\;840,\;720
\]
are checked by kernel reduction, the top one through
$\mathrm{surj}(n)=n!$, and the target \eqref{eq:target} follows by one
\texttt{field\_simp}.  The counts coincide with Atlas 148's, as they must,
but they are checked again: the two graphs are not isomorphic, so nothing
transports.

\subsection{Scope}
\label{sec:app145-scope}

The formalization asserts exactly the corollary: the two stated page
placements on a $4$-cycle, in that direction.  It does not formalize the
general relocation principle discussed in Section~4, and --- as that section
records --- the opposite-edge placement is Atlas 152, which has a verified
counterexample in this project
(\texttt{Methods.Negative.Atlas152.violatesLowerBound\_152}), so no
one-sided concentration inequality could cover it.

With this row the catalogue reaches $93$ verified rows: every Atlas graph
whose status is currently accepted, positive or negative, is proved in Lean
with no \texttt{sorry}.

\end{document}
